%!TEX program = xelatex
\documentclass{md2tex}

% 参考文献设置（仅正文存在 citation 时保留）
$if(bibliography)$
\usepackage{gbt7714}
\citestyle{numbers}
\renewcommand{\bibname}{参考文献}
$endif$

% 论文封面信息
$if(title)$
\newcommand{\papertitle}{$title$}
$else$
\newcommand{\papertitle}{文档标题}
$endif$
\newcommand{\securitylevel}{$security$}
% 保密年限：与密级分列两个变量，非空时以五角星（pifont 的 \ding{72}）连接
\newcommand{\securityyears}{$securityyears$}
\newcommand{\securitymark}{%
  \securitylevel\ifdefempty{\securityyears}{}{\ding{72}\securityyears}}
\newcommand{\doctype}{$doctype$}
\newcommand{\docnumber}{$docnumber$}
\newcommand{\docversion}{$docversion$}
\newcommand{\institution}{$institution$}
\newcommand{\dateofsubmit}{$submitdate$}
% 封面专用：公开件不标密级；稿次已带全角括号；阶段 0–3 为项目类，-1 为研究类
\newcommand{\coversecurity}{$coversecurity$}
\newcommand{\coverversion}{$versionmark$}
\newcommand{\coverident}{$ident$}
\newcommand{\coverbyline}{$byline$}
\newcommand{\coveroriginal}{$originaltitle$}
\newcommand{\coverstage}{$coverstage$}
\newcommand{\coverdoctype}{$doctypespaced$}
% 封面题名：程序按词边界算好分行，行间以 \\ 断开
\newcommand{\covertitle}{$covertitle$}

\begin{document}

% 封面：位置与字号见 mdx `cover::layout`，Word 与宿主预览按同一组数排。
% 整页用 TikZ 按距页顶的毫米数定位，不随题名行数上下漂移。
% 文种的字距由程序逐字插入 \hspace，不用 \ziju：它会漏到后面的题名节点里。
\begin{titlepage}
  \thispagestyle{empty}
  \null
  \begin{tikzpicture}[remember picture, overlay,
      every node/.style={inner sep=0pt, outer sep=0pt}]
    \coordinate (P) at (current page.north west);
    % 密级（左）与编号（右）
    \ifdefempty{\coversecurity}{}{%
      \node[anchor=north west, font=\heiti\enhei\zihao{-4}]
        at ([shift={(25mm,-20mm)}]P)
        {\coversecurity\ifdefempty{\securityyears}{}{\ding{72}\securityyears}};}
    \ifdefempty{\docnumber}{}{%
      \node[anchor=north east, font=\heiti\enhei\zihao{-4}]
        at ([shift={(185mm,-20mm)}]P) {编号：\docnumber};}
    % 文种：黑体小二，字距半字，黑色
    \node[anchor=north, font=\heiti\enhei\zihao{-2}]
      at ([shift={(105mm,-68mm)}]P) {\coverdoctype};
    % 标识行
    \ifdefempty{\coverident}{}{%
      \node[anchor=north, font=\song\zihao{4}]
        at ([shift={(105mm,-83mm)}]P) {\coverident};}
    % 反线：项目类单线，研究类一粗一细
    \fill ([shift={(25mm,-94mm)}]P) rectangle ([shift={(185mm,-94.5mm)}]P);
    \ifnum\coverstage<0
      \fill ([shift={(25mm,-95.3mm)}]P) rectangle ([shift={(185mm,-95.5mm)}]P);
    \fi
    % 题名：小标宋一号，行距 1.45 倍
    \node[anchor=north, text width=150mm, align=flush center]
      (T) at ([shift={(105mm,-114mm)}]P)
      {\bs\fontsize{26bp}{37.7bp}\selectfont\covertitle\par};
    \coordinate (L) at (T.south);
    \ifdefempty{\coverversion}{}{%
      \node[anchor=north, font=\song\zihao{-3}]
        (V) at ([yshift=-6mm]L) {\coverversion};
      \coordinate (L) at (V.south);}
    \ifdefempty{\coveroriginal}{}{%
      \node[anchor=north, text width=150mm, align=flush center]
        at ([yshift=-6mm]L)
        {\itshape\fontsize{15bp}{20bp}\selectfont\coveroriginal\par};}
    % 项目类：四格阶段条，当前阶段黑色粗线，其余灰色细线
    \ifnum\coverstage>-1
      \foreach \i/\name in {0/立项论证,1/建设实施,2/技术实现,3/项目总结} {
        \pgfmathsetmacro{\xl}{25+\i*40.5}
        \pgfmathsetmacro{\xr}{\xl+38.5}
        \pgfmathsetmacro{\xc}{\xl+19.25}
        \ifnum\i=\coverstage
          \fill ([shift={(\xl mm,-218mm)}]P) rectangle ([shift={(\xr mm,-218.8mm)}]P);
          \node[anchor=north, font=\heiti\enhei\zihao{5}]
            at ([shift={(\xc mm,-221.2mm)}]P) {\name};
        \else
          \fill[black!35] ([shift={(\xl mm,-218mm)}]P) rectangle ([shift={(\xr mm,-218.2mm)}]P);
          \node[anchor=north, font=\heiti\enhei\zihao{5}, text=black!45]
            at ([shift={(\xc mm,-221.2mm)}]P) {\name};
        \fi
      }
    \fi
    % 研究类：署名行
    \ifnum\coverstage<0
      \ifdefempty{\coverbyline}{}{%
        \node[anchor=north, font=\song\zihao{4}]
          at ([shift={(105mm,-224mm)}]P) {\coverbyline};}
    \fi
    % 落款：单位黑体三号，日期宋体小三
    \node[anchor=north, font=\heiti\enhei\zihao{3}]
      (O) at ([shift={(105mm,-245mm)}]P) {\institution};
    \node[anchor=north, font=\song\zihao{-3}]
      at ([yshift=-5mm]O.south) {\dateofsubmit};
  \end{tikzpicture}
\end{titlepage}
\clearemptydoublepage

% 正文从阿拉伯页码 1 起。目录不在这里排：正文写了 `<!-- [目录] -->` 才在
% 标记所在处排（见 md2tex.cls 的目录命令，单用大写罗马页码）
\mainmatter

\vspace{28pt}

% 正文
$body$

% 参考文献
$if(bibliography)$
\bibliography{references}
$endif$

\end{document}
